\documentclass[aspectratio=169]{beamer}
\input{header}
\coursecode{JKSSB}

\title{Computer Generations \& History}
\subtitle{Lesson 49 --- Technology Milestones and Association Traps}
\author{JKSSB Computer Awareness}
\date{\today}

\begin{document}
\frame{\titlepage}

\begin{frame}{What Does a Generation Label Mean?}
  \begin{itemize}
    \item Textbooks group computer history by the dominant hardware
      technology.
    \item A generation label is a learning classification, not a precise
      date boundary used identically by every source.
    \item For exams, learn the technology association first.
  \end{itemize}
  \begin{block}{High-yield distinction}
    Integrated circuits (ICs) are associated with the third generation;
    microprocessors with the fourth.
  \end{block}
\end{frame}

\begin{frame}{First and Second Generations}
  \begin{columns}[T]
    \column{0.46\textwidth}
      \textbf{First generation}
      \begin{itemize}
        \item Vacuum tubes
        \item Large, heat-producing systems
        \item Machine-language programming
      \end{itemize}
    \column{0.46\textwidth}
      \textbf{Second generation}
      \begin{itemize}
        \item Transistors replace vacuum tubes
        \item Smaller and more reliable systems
        \item Assembly and early high-level languages
      \end{itemize}
  \end{columns}
\end{frame}

\begin{frame}{Third and Fourth Generations}
  \begin{columns}[T]
    \column{0.46\textwidth}
      \textbf{Third generation}
      \begin{itemize}
        \item Integrated circuits (ICs)
        \item Greater reliability and reduced size
        \item Operating systems and multiprogramming develop
      \end{itemize}
    \column{0.46\textwidth}
      \textbf{Fourth generation}
      \begin{itemize}
        \item Microprocessors
        \item VLSI integration
        \item Personal computers become widespread
      \end{itemize}
  \end{columns}
\end{frame}

\begin{frame}{Fifth Generation: Use Qualified Wording}
  \begin{itemize}
    \item Introductory texts often associate the fifth generation with
      Artificial Intelligence (AI), parallel processing, and natural-language
      systems.
    \item This is a broad textbook grouping, not a universally fixed
      technology boundary.
    \item Do not confuse an AI application with the microprocessor milestone
      used for the fourth generation.
  \end{itemize}
  \begin{alertblock}{Exam method}
    Match the wording asked in the question to the conventional association;
    avoid claiming that every modern computer belongs to one sharply dated era.
  \end{alertblock}
\end{frame}

\begin{frame}{India Milestone: C-DAC and PARAM 8000}
  \begin{itemize}
    \item \key{C-DAC}: Centre for Development of Advanced Computing.
    \item \key{PARAM 8000}: India's first supercomputer, developed by C-DAC
      and unveiled in 1991.
    \item It belongs to India's high-performance computing history, not to
      the definition of a computer generation.
  \end{itemize}
  \vspace{0.3cm}
  \href{https://cdac.in/index.aspx?id=hpc_ss_supercomputing_systems}
       {C-DAC: HPC Supercomputing Systems}
\end{frame}

\begin{frame}{Association Ladder}
  \begin{center}
    \begin{tabular}{ll}
      \toprule
      \textbf{Milestone} & \textbf{Association} \\
      \midrule
      Vacuum tubes & First generation \\
      Transistors & Second generation \\
      Integrated circuits & Third generation \\
      Microprocessors & Fourth generation \\
      AI / parallel processing (conventional grouping) & Fifth generation \\
      \bottomrule
    \end{tabular}
  \end{center}
  \vspace{0.3cm}
  \muted{Anchor the sequence: tubes, transistors, ICs, microprocessors.}
\end{frame}

\begin{frame}{Lesson 49: Recall}
  \begin{itemize}
    \item Third generation = ICs; fourth generation = microprocessors.
    \item First = vacuum tubes; second = transistors.
    \item Fifth-generation descriptions commonly mention AI, but boundaries
      are conventional.
    \item C-DAC developed PARAM 8000, India's first supercomputer, unveiled
      in 1991.
  \end{itemize}
\end{frame}
\end{document}
